
%\needspace{4\baselineskip}
\color{uniblue}{\subsection{Междуфазная токовая отсечка}}
\color{black}

\begin{enumerate}[label=\arabic{section}.\arabic{subsection}.\arabic{enumi}, labelsep=4pt, leftmargin=0pt, itemindent=57pt, itemsep=0pt, parsep=5pt] % если что можно крутить itemsep
%%%%%%%%%%%%%%%%%%%%%%%%%%%%%%%%%%%%%%%%%%%%%%%%%%%%%%%%%%%%%%%%%%%%%%  ОБЩИЕ СВЕДЕНИЯ %%%%%%%%%%%%%%%%%%%%%%%%%%%%%%%%%%%%%%%%%%%%%%%%%%

\item
Для резервирования действия защит при близких КЗ в устройстве предусмотрена ненаправленная 
междуфазная токовая отсечка (МФТО), работающая по факту превышения уставок не менее чем двумя из трех 
фазных токов

\begin{figure}[H]
\centering
\includegraphics[width=0.95\textwidth,height=0.95\textheight,keepaspectratio]{\fbpath/561.0171 МФТО/img/MFTO.pdf}
\captionof{figure}{Функциональная схема МФТО}\label{pict_bk:MFTO}
\end{figure}

\item
Ввод в работу МФТО осуществляется программным ключом \textbf{«МФТО»}.

\item
Ток срабатывания МФТО задается с помощью уставки \textbf{«Iср МФТО»}. Время срабатывания задаётся 
уставкой \textbf{«Тср МФТО»}. 

\item
Для предотвращения неправильной работы МФТО предусмотрена блокировка при протекании через 
трансформаторы тока двух выключателей сквозного тока внешнего КЗ, вызывающего неодновременное 
насыщение трансформаторов. Блокировка вводится программным ключом \textbf{«БСТО МФТО»}.

\item
МФТО может быть оперативно выведена из работы от внешнего сигнала \textit{«Вывод МФТО»}, назначенного
на соответствующий дискретный вход или функциональную клавишу.








\end{enumerate}
